% !TEX root = ../main.tex
\chapter{Conclusion and Future Work}
\label{ch:conclusion}

Chapter~\ref{ch:discussion} interpreted the results and judged the objectives met. This chapter closes the thesis: it answers each research question in one paragraph, states the original contribution to knowledge without repeating the contribution list of Chapter~\ref{ch:introduction}, and sets out the work that would relax the limits recorded in Section~\ref{sec:limitations}. Numbers are not restated; the tables of Chapter~\ref{ch:results} remain the single point of reference.

\dissertationSubheading[sec:summary-findings]{Answers to the research questions}

The thesis formulated EndorseRank, PageRank on a latest-allowance endorsement graph, and evaluated it against Adaptive Weighted PageRank (AWP), the transfer-graph baseline of Do, Do and Nguyen\cend{3}, on one matched Arbitrum One wallet cohort, one spender holdout cohort and one expanded address pool over December~2025 to May~2026.

RQ1 asked whether EndorseRank and AWP produce meaningfully different wallet orderings on an identical sample. They do. The inter-method rank agreement lies strictly between unrelated and identical orderings with an interval that excludes both extremes (Table~\ref{tab:alignment-family-ci}), and the rank scatter shows substantial pairwise disagreement. Hypothesis~H1 is supported: authorisation graphs and payment graphs are not interchangeable reputation substrates.

RQ2 asked which method aligns more strongly with contemporaneous allowance checks and with contemporaneous transfer checks. The answer is construct-specific. AWP aligns more strongly with transfer proxies and EndorseRank with allowance proxies, and the two between-method paired differences exclude zero in opposite directions (Table~\ref{tab:tau-diff}). H2 is supported in that form. The within-method differences---EndorseRank's allowance alignment against its own transfer alignment---did not reach significance, so the thesis does not claim that EndorseRank favours one signal over the other inside its own ranking.

RQ3 asked whether EndorseRank is computationally more efficient than AWP on the same matched sample. It is, on wall-clock time, memory and iterations, and the advantage is attributable to the endorsement graph having about one ninth of the transfer graph's edges under a single timing protocol (Table~\ref{tab:benchmark-runtime}). H3 is supported. The node-count sensitivity on the expanded pool confirms the ratio across staged node counts but, because the second-tier extraction did not complete, does not test edge growth (Table~\ref{tab:benchmark-scaling}).

RQ4 asked whether scores frozen at t1 predict new owner--spender approvals at t2 on inbound spenders more strongly than AWP and more informatively than the same-window allowance check. They do (Table~\ref{tab:holdout-spenders}): EndorseRank's interval excludes zero and lies above AWP's, whose interval includes zero, and the labels post-date every scored event. RQ4 replaced the trading-success comparison of the research proposal, which was withdrawn as circular (Section~\ref{sub:gf-pr}); the holdout tests temporal predictive validity of the endorsement construct, not credit-risk validity.

Taken together, Claims~C1--C4 (Section~\ref{sec:empirical-claims}) establish an efficiency--alignment trade-off rather than a dominance: EndorseRank offers lower PageRank cost, allowance-construct alignment and out-of-window evidence on future approvals; AWP offers stronger transfer and longitudinal-stability alignment. Which to deploy depends on whether the operator's question is about authorisation or about flow.

\dissertationSubheading[sec:contribution-knowledge]{Original contribution to knowledge}

The contribution list of Section~\ref{sec:contributions} named what the thesis set out to add. What it has in fact established, in the terms of the literature reviewed in Chapter~\ref{ch:literature}, is the following.

Before this work, the graph-reputation literature on public ledgers ranked wallets on transfers \cend{4,3}, on transaction-network risk propagation \cend{12} or on explicit ratings and stake \cend{25,26}, and it validated those rankings against statistics drawn from the same input data. The ERC-20 allowance---the one on-chain act that records a deliberate delegation of spending authority \cend{18}---had not been used as a reputation edge, and no same-cohort comparison with confidence intervals, difference tests and a temporal holdout existed for any pair of on-chain reputation methods. This thesis supplies that edge construct and that comparison. Its knowledge claim is not that a new algorithm outperforms an old one; PageRank is used unchanged \cend{15}. It is that the choice of edge determines what an on-chain reputation score measures, that this can be shown with quantified uncertainty on one cohort, and that the allowance edge yields a score which is cheaper to refresh, distinct from the transfer-based score, aligned with its own construct and predictive of later endorsement---while being weaker than the transfer-based score on flow and stability. The construct-validity framing \cend{22,33} that makes this a testable statement, and the open pipeline that makes it a reproducible one, are part of the contribution, because the auditability of decentralised credit-related scores is itself a recognised concern \cend{2}.

The research is deliberately bounded. It does not show that on-chain reputation can substitute for collateral \cend{9,1}, that EndorseRank is Sybil-proof (Section~\ref{sec:sybil-cost-model}) or that its results hold beyond Arbitrum One and the studied window.

\dissertationSubheading[sec:future-work]{Future work}

Each item below is tied to a limitation recorded in Chapter~\ref{ch:discussion}.

\begin{itemize}
\item Completing the second-tier extraction. The expanded pool tested node count only. Extracting approval and transfer logs for the supplemental addresses would allow runtime to be measured as the edge population grows and would turn the bounded node-count statement of RQ3 into a scaling result.
\item Cross-protocol and cross-chain validation. Replication on other rollups and Layer-1 networks would test whether the allowance--transfer trade-off and the holdout pattern generalise beyond Arbitrum One; the heterogeneity of DeFi venues \cend{8,1} makes this the main external-validity gap.
\item Longer windows and a credit-linked holdout. A longer observation window would test stability of the coefficients across market regimes, and a holdout label tied to repayment or liquidation, where lending-protocol coverage of the spender cohort permits, would test whether the endorsement construct carries any credit-risk information; the present holdout does not.
\item Adversarial evaluation. The Sybil cost model is analytical. Empirical red-teaming---wash approvals, coordinated spender farms, strategic revocation---would quantify realised attack cost-effectiveness under live fee markets and detection heuristics \cend{11}, and would test whether the Sybil-adjusted stability proxies remain informative under active adversaries.
\item Hybrid and richer authorisation semantics. Combining endorsement-graph and transfer-graph features, or adding protocol-specific delegation logs to the edge definition, may improve coverage; such designs should report construct-specific checks and time-split evidence rather than a single aggregate metric, and should preserve the auditability that distinguishes graph propagation from opaque learned scores \cend{14,2}.
\item Deployment studies. Integration with risk oracles and governance tooling would test whether EndorseRank's runtime profile sustains production refresh cadences, and controlled fee-tier or monitoring-cadence experiments would replace the adoption hypotheses of Section~\ref{sub:deploy-industry} with evidence. Collateral reduction remains a hypothesis, not a claim of this research.
\end{itemize}

\dissertationSubheading[sec:closing-reflection]{Closing reflection}

On-chain reputation will not dissolve the need for collateral in adversarial, pseudonymous markets \cend{9}. It can make behavioural structure legible, and this research shows that legibility depends on which edges one chooses to read. Transfers tell a story about economic flow; allowances tell a story about delegated authority. PageRank can narrate either, but not both at once with the same ranking. EndorseRank commits to the authorisation narrative and documents its costs and benefits, including a time-split check of later approvals and two honest null results about its own internal structure.

The practical upshot is modest and actionable. A protocol or researcher who needs a fast, auditable score of who is trusted to spend has, on Arbitrum One data of the kind studied here, a credible instrument in EndorseRank. One who needs to find transfer hubs should still use AWP. One who wants to predict trading profit from social graphs alone is asking a question this thesis has shown to be outside what either instrument can answer.
